%! TeX program = lualatex
\documentclass[../main.tex]{subfiles}

\usetikzlibrary{intersections,calc}

\newenvironment{solution}{
  \begin{proof}[Sample solution]
    \color{main}
  }{
  \end{proof}
}

\begin{document}

\chapter*{Week 11 practice problems}

\setcounter{chapter}{11}

\begin{exercise}
  Suppose \(A = \begin{bmatrix} 0 & 1 \\ 1 & 0\end{bmatrix}\). Calculate \(A^{2}, A^{3}, A^{4}\). Find a pattern to calculate \(A^{53}\).
\end{exercise}

\begin{exercise}
  Suppose \(A = \begin{bmatrix} 1 & 1 \\ 0 & 1\end{bmatrix}\). Calculate \(A^{2}, A^{3}, A^{4}\). Find a pattern to calculate \(A^{2026}\).
\end{exercise}

\begin{exercise}
  Suppose \(A = \begin{bmatrix} 12 & -3  \\ 8 & 2 \end{bmatrix}\) is a matrix.  It is known that \(\vec{u} = \begin{bmatrix} 1 \\ 2 \end{bmatrix}\) and \(\vec{v} = \begin{bmatrix} 3 \\ 4 \end{bmatrix}\) are its eigenvectors.

  \begin{enumerate}
    \item Calculate \(A^{3}\vec{u}\) by brute-force. What is the eigenvalue for \(\vec{u}\)? How many powers of \(A^{t}\vec{u}\) do you need to calculate before you can recognize this eigenvalue?
    \item Find the eigenvalue for \(\vec{v}\).
    \item Suppose \(\vec{x}(0) = \begin{bmatrix} 2 \\ 3 \end{bmatrix}\) and \(\vec{x}(t+1) = A \vec{x}(t)\) for all \(t \ge 1\). Find \(\vec{x}(1000)\).
  \end{enumerate}
\end{exercise}

\begin{exercise}
  A matrix \(A\) has an eigenvector \(\vec{v} = \begin{bmatrix}1 \\ 3 \\ -7\end{bmatrix}\). We know that \(A^{2} \vec{v} = \begin{bmatrix} 25 \\ 75 \\ -175\end{bmatrix}\). What is the eigenvalue for \(\vec{v}\)?
\end{exercise}

\begin{exercise}
  Find all eigenvectors and their associated eigenvalues for \(A = \begin{bmatrix} \sqrt{3} & -1 \\ 1 & \sqrt{3} \end{bmatrix}\).
\end{exercise}

\begin{exercise}
  Find all eigenvectors and their associated eigenvalues for \(A = \begin{bmatrix} 3 & 0 \\ 0 & -2 \end{bmatrix}\).
\end{exercise}

\begin{exercise}
  Find all eigenvectors and their associated eigenvalues for \(A = \begin{bmatrix} 0 & -1 \\ 2 & 1 \end{bmatrix}\).
\end{exercise}

\begin{exercise}
  Find all eigenvectors and their associated eigenvalues for \(A = \begin{bmatrix} 2 & -2 \\ 1 & -3 \end{bmatrix}\).
\end{exercise}

\begin{exercise}
  Find all eigenvectors and their associated eigenvalues for \(A = \begin{bmatrix} 2 & 1 \\ 2 & 3 \end{bmatrix}\).
\end{exercise}

\begin{exercise}
  Find all eigenvectors and their associated eigenvalues for \(A = \begin{bmatrix} 2 & 2 \\ 0 & 3 \end{bmatrix}\).
\end{exercise}

\begin{exercise}
  Find all eigenvectors and their associated eigenvalues for \(A = \begin{bmatrix} 0 & -1 \\ 1 & 2 \end{bmatrix}\).
  
  % \begin{solution}
  %   There is exactly one eigenvalue \(\lambda = 1\) with eigenvectors \(\begin{bmatrix} -1 \\ 1 \end{bmatrix}\).
  % \end{solution}
\end{exercise}

\begin{exercise}
  Find all eigenvectors and their associated eigenvalues for \(A = \begin{bmatrix} 1/2 & 1 \\ -1 & 1/3 \end{bmatrix}\).
\end{exercise}

\begin{exercise}
  Find all eigenvectors and their associated eigenvalues for \(A = \begin{bmatrix} 3 & 0 & 0 \\ 1 & 2 & 0 \\ 0 & 1 & 1 \end{bmatrix}\).
\end{exercise}

\begin{exercise}
  Find all eigenvectors and their associated eigenvalues for \(A = \begin{bmatrix} 1 & 2 & 0 \\ 3 & 4 & 0 \\ 5 & 6 & 7 \end{bmatrix}\).
\end{exercise}
\end{document}
